% TOP_PROBLEM: {"id": 1110, "title": "Herzog-Schönheim Conjecture", "category_id": 4, "category": {"id": 4, "name": "algebra", "display_name": "Algebra", "description": "Group theory, ring theory, field theory, and algebraic structures.", "slug": "algebra", "order_index": 4, "created_at": "2026-07-31T15:26:25.671Z"}, "status": "open"}
% FIRST_POSTED: null
% ENTRY_KIND: statement_only
% Imported from: batch04.tex; UnsolvedMath ID 1110
\documentclass[11pt]{article}
\usepackage[margin=25mm]{geometry}
\usepackage{fontspec}
\setmainfont{Latin Modern Roman}
\usepackage{amsmath,amssymb,mathtools}
\usepackage{unicode-math}
\setmathfont{Latin Modern Math}
\usepackage{xurl,hyperref}
\hypersetup{colorlinks=true,urlcolor=blue}
\setcounter{secnumdepth}{0}
\setlength{\parindent}{0pt}
\setlength{\parskip}{5pt}
\setlength{\emergencystretch}{3em}
\newcommand{\sref}[2]{\hyperlink{source:#1:#2}{[S#2]}}
\newcommand{\cataloguescope}[1]{\paragraph{Recorded target.}#1}
\begin{document}
% UnsolvedMath ID 1110; source catalogue code ALG-010
\setcounter{section}{1109}
\section{Herzog-Schönheim Conjecture}\label{1110}
{\small\sffamily UnsolvedMath ID 1110 · Algebra}
\subsection{Definitions and mathematical statement}

\paragraph{Shared notation.}$\mathbb N=\{1,2,\ldots\}$, and $\mathbb Z,\mathbb Q,\mathbb R,\mathbb C$ denote the integers, rationals, reals and complex numbers. Any other conventions are specified locally.


For a subgroup \(H\le G\), the index \([G:H]\) is the number of its
left cosets.  A nontrivial finite coset partition is a disjoint union
\[
 G=\bigsqcup_{i=1}^{s} g_iH_i,\qquad s\ge2,
\]
where each \(g_i\in G\) and each \(H_i<G\) is a subgroup of finite index.
There is no assumption that \(G\) is finite, abelian or finitely generated.
The nontriviality requirement excludes the one-part partition \(G=G\).
Right-coset and left-coset formulations are equivalent by inversion.
\paragraph{Statement recorded in the supplied source.}
For every group \(G\) and every nontrivial finite coset partition
\(G=\bigsqcup_{i=1}^{s}g_iH_i\), must there be distinct indices
\(i,j\in\{1,\ldots,s\}\) such that
\[
 [G:H_i]=[G:H_j]?
\] \sref{1110}{2}
\subsection{Short English statement}
In every nontrivial partition of a group into finitely many subgroup cosets, must two of the subgroups have the same index?
\subsection{Sources}
\begin{description}
\item[\hypertarget{source:1110:1}{S1}] ULAM.AI and the UnsolvedMath Contributors, \emph{UnsolvedMath: A Curated Collection of Open Mathematics Problems}, record 1110 (ALG-010). CC BY 4.0. \url{https://huggingface.co/datasets/ulamai/UnsolvedMath}.
\item[\hypertarget{source:1110:2}{S2}] Fabienne Chouraqui, \emph{The Herzog--Sch\"onheim conjecture for finitely generated groups}, arXiv:1803.08301 (2018), Section 1, including the formulation for arbitrary groups. \url{https://arxiv.org/abs/1803.08301}.
\end{description}
\label{end:1110}
\end{document}
